\hypertarget{appendix-0fc95374cb7e}{}
\section*{Render graph barrier lowering}
\section{Render graph barrier lowering}\label{render-graph-barrier-lowering}

A logical render graph and an API synchronization plan are different layers.

The logical graph answers:

\begin{itemize}
\tightlist
\item
  which pass reads or writes each resource;
\item
  which hazards create ordering edges;
\item
  whether the pass dependency graph is acyclic;
\item
  when a logical resource is first and last used.
\end{itemize}

API lowering answers:

\begin{itemize}
\tightlist
\item
  what concrete stage/access state each use requires;
\item
  what image layout or resource state is required;
\item
  which subresources are affected;
\item
  which queue executes each pass;
\item
  what barrier, semaphore, fence/event, or encoder boundary realizes the dependency.
\end{itemize}

The repository already has tools/labs/render\_graph\_contract\_lab.cpp for the first layer. That CPU witness proves RAW/WAR/WAW edge construction, deterministic topological order and several invalid-graph cases. It does not prove Vulkan or Metal memory visibility.

\subsection{State machine per resource version}\label{state-machine-per-resource-version}

A practical compiler treats each write as producing a new logical version. A read consumes a particular version. For every concrete resource/subresource, track the last known backend state:

producer pass\\
queue\\
stage/access\\
layout or storage state\\
version.

When the next use requires an incompatible state, emit a transition whose source scope covers the prior access and whose destination scope covers the next access.

This versioned model prevents an important bug: a pass can read ``history texture'' while the graph accidentally associates it with the current-frame version rather than the previous-frame version.

\subsection{Subresource granularity}\label{subresource-granularity}

Whole-resource tracking is simple but can over-synchronize. Independent mip levels, array layers or depth/stencil aspects may be used by different passes. Moving to subresource tracking requires interval/range splitting and merging so that state for mip 0 does not overwrite the state of mip 5.

The first correct implementation should prefer conservative whole-resource barriers over an incorrect fine-grained optimizer. Fine granularity is an optimization after invariants are tested.

\subsection{Vulkan lowering}\label{vulkan-lowering}

For Vulkan synchronization2, a transition record eventually becomes VkBufferMemoryBarrier2, VkImageMemoryBarrier2 or VkMemoryBarrier2 inside VkDependencyInfo, possibly combined with queue submission dependencies.

See \hyperref[appendix-f15ce0d9c170]{Vulkan synchronization2 and resource transitions}.

\subsection{Metal lowering}\label{metal-lowering}

Metal does not mirror Vulkan's barrier vocabulary one-for-one. The backend should lower the same logical dependency into Metal's encoder ordering, resource hazard mode, fences/events where required, and heap/lifetime rules rather than forcing Metal to imitate Vulkan object names.

See \hyperref[appendix-563b2cb23fc3]{Metal heaps memoryless resources and hazards}.

\subsection{Verification ladder}\label{verification-ladder}

A render-graph compiler should be verified in layers:

\begin{enumerate}
\def\labelenumi{\arabic{enumi}.}
\tightlist
\item
  dependency graph against a CPU oracle;
\item
  resource-state transition plan against hand-written traces;
\item
  Vulkan execution with validation enabled;
\item
  Metal execution with GPU capture;
\item
  cross-backend image/data parity;
\item
  transient aliasing stress with deliberately overlapping and non-overlapping lifetimes;
\item
  current production-engine integration.
\end{enumerate}

The current branch now has stages 1 and a minimal Vulkan synchronization2 execution/Metal compile boundary. The deeper stages remain explicit rather than implied.
